\chapter{Planificación y presupuesto}
\label{cap:planificacion}

\section{Fases}
\label{sec:fases}
La planificación del proyecto se estructurará en fases con el objetivo de garantizar un desarrollo organizado y eficiente, asegurando que cada etapa esté enfocada en el cumplimiento de objetivos específicos.

\subsection{Fase Inicial}
En esta fase se recopilará y definirá toda la información relacionada con el Capitulo 1, incluyendo la motivación para la realización del proyecto y la determinación de los objetivos que se pretenden alcanzar a lo largo del mismo. 

\subsection{Fase de Investigación}
Esta fase estará dedicada a la recopilación de información relevante para el desarrollo del proyecto. Se analizarán aplicaciones similares para demostrar la necesidad de una aplicación integral que cubra todas las funciones propuestas. 
Además se investigará aspectos como la gestión del tiempo en hogares compartidos y la administración de gastos entre personas que comparten un mismo domicilio.

Para garantizar que la aplicación cumpla con los requisitos legales, también se revisará la normativa vigente de protección de datos. Asímismo, se realizará un estudio 
de usabilidad y accesibilidad para asegurar que el diseño sea inclusivo y fácil de usar para todos los usuarios.

\subsection{Fase de Analisis}
En esta fase se llevará a cabo un análisis de las diferentes tecnologías que se pueden utilizar para el desarrollo de la aplicación, tanto para el frontend como para el 
backend, con el objetivo de selecionar las más adecuadas para el proyecto. Además, se definirán los requisitos que debe cumplir la aplicación para satisfacer las necesidades del usuario objetivo. Como
parte de esta fase, se elaborarán los casos de uso y los diagramas de secuencia del sistema.

\subsection{Fase de Diseño}
Esta fase estará enfocada en la creación de los prototipos de la aplicación. Se desarrollará un diagrama de arquitectura del sistema y bocetos para la interfaz gráfica. Además, se definirán 
los colores y las tipografías que se emplearán en la aplicación, garantizando un diseño responsive que permita su visualización en cualquier dispositivo.

\subsection{Fase de Desarrollo}
En esta fase se procederá a la implementación de la aplicación, siguiendo las tecnologías, metodologías y diseños definidos en las fases previas de análisis y diseño. Esta fase se dividirá en cuatro iteraciones,
ya que se empleará una metodología ágil para el desarrollo del proyecto. En cada iteración se determinarán las funcionalidades que se implementarán y se realizarán pruebas para garantizar 
que la aplicación cumple con los requisitos definidos en la fase de análisis.

Las pruebas serán realizadas por usuarios reales, concretamente jóvenes que comparten piso, para verificar que la aplicación satisface sus necesidades y resulta fácil de utilizar.

\subsection{Documentación}
Aunque no se considera una fase independiente, la documentación será un proceso continuo a lo largo de todo el proyecto. Se registrarán de manera detallada todas las actividades y resultados obtenidos en cada fase. 

\subsection{Diagrama de Gantt}
En esta tabla se específica la duración en horas de cada una de las tareas representadas en el diagrama de Gantt. Este diagrama representa 
la planificación temporal de las actividades del proyecto, donde cada celda corresponde a una semana del mes. La tabla presenta la cantidad de horas
dedicadas a cada tarea, mientras que el diagrama de Gantt muestra el periodo en el que se llevaron a cabo. 

\begin{table}[ht]
    \centering
    \resizebox{\textwidth}{!}{%
    \begin{tabular}{|l|c|}
        \hline
        \rowcolor[HTML]{C0C0C0} 
        \textbf{Tarea} & \textbf{Horas} \\ \hline
        Motivación & 1 \\ \hline
        Objetivos del proyecto & 2\\ \hline
        Estructura de la memoria  & 2 \\ \hline
        Planificación & 3 \\ \hline
        Presupuesto & 4\\ \hline
        Estudios sobre la gestión del hogar y las finanzas en viviendas compartidas & 20 \\ \hline
        Aplicaciones similares y comparativas  & 25 \\ \hline
        Legislación actual de protección de datos & 7 \\ \hline
        Estudio de usabilidad y accesibilidad & 12 \\ \hline
        Análisis de las tecnologías Frontend & 20 \\ \hline
        Analisis de las tecnologías Backend & 15 \\ \hline
        Análisis de las metodologías de desarrollo & 12 \\ \hline
        Explicación de la propuesta & 2 \\ \hline
        Análisis de la aplicación & 40 \\ \hline
        Diseño de la aplicación & 40 \\ \hline
        Elección de tecnologias & 2 \\ \hline
        Iteración 1 & 30 \\ \hline
        Iteración 2 & 45 \\ \hline
        Iteración 3 & 45 \\ \hline
    \end{tabular} %
    }
    \caption{Tabla de tareas y horas}

\end{table}

\newpage

\section{Presupuesto}
\label{sec:presupuesto}
En este apartado se detallarán los costes asociados al desarrollo del proyecto. Se incluirán los costes de hardware, software, licencias, servicios en la nube y otros gastos necesarios para la realización del proyecto.
Realizar un presupuesto en este proyecto es clave para planificar y gestionar los recursos de forma eficiente. Permite asegurarse de que las fases de desarrollo se puedan llevar a cabo con los medios disponibles. 
Además, ayuda a justificar las decisiones tomadas y evitar problemas por faltas de recursos asegurando que el proyecto sea viable y esté bien estructurado. 

\subsection{Recursos Hardware}
Los componentes hardware necesarios para el desarrollo del proyecto son un ordenador portátil y un teléfono móvil para realizar pruebas con la aplicación. 

\begin{itemize}
    \item \textbf{Ordenador:} TUF Gaming FX505DT\_FX505DT 
    \item \textbf{Teléfono Móvil:} Samsung Galaxy A51
\end{itemize}

\subsection{Recursos Software}
Las herramientas software que se utilizarán en el proyecto son las siguientes:
\begin{itemize}
    \item \textbf{Sistema Operativo:} Windows 11
    \item \textbf{IDE:} Visual Studio Code
    \item \textbf{Control de Versiones:} Git
    \item \textbf{Lenguajes de Programación:} Flutter, Dart
    \item \textbf{Diseño de diagramas UML:} Visual Paradigm
    \item \textbf{Diseño de Interfaces:} Figma
    \item \textbf{Sistema de Gestión de Base de Datos:} Firebase
    \item \textbf{Sistema de composición de texto:} LaTeX
\end{itemize}

\subsubsection{Recursos Humanos}
En este apartado se detallarán los gastos por la contratación de personal. 
Este proyecto será desarrollado por una única persona, un programador junior, que según \cite{Talent2025} tiene un coste de 21.000\euro al año, como la duración de este proyecto es de 6 meses, el coste total será de 10.500\euro.


\subsection{Licencias}
En este apartado se detallan las licencias software necesarias para llevar a cabo el desarrollo del proyecto. Aunque en este caso no supone ningún coste, 
ya que se utilizarán versiones gratuitas o herramientas sin coste, en un proyecto real sería imprescindible considerar el pago de las licencias. Esto garantizaría 
el acceso completo a todas las funcionalidades de las herramientas, además de cumplir con las normativas legales y éticas relaciondas con el uso del sotware. Las licencias de software que deberían considerarse en un proyecto real 
son las siguientes: 
\begin{itemize}
    \item \textbf{Windows 11:} Windows 11 Home License
    \item \textbf{Visual Paradigm:} Visual Paradigm Standard Edition
    \item \textbf{Figma:} Figma Professional
\end{itemize}

El resto de licencias necesarias para el desarrollo del proyecto son gratuitas, por lo que no se incluirán en el presupuesto.


\subsection{Costes indirectos}
En este apartado se engloba todos los gastos indirectos que incluyen los costes de electricidad, internet, espacio de trabajo, etc. 
Estos gastos indirectos suelen estimarse como un porcentaje adicional del coste directo total \textit{(entre 10\% y 20\%)}.
En este caso será un 15\% de los recursos humanos.

\subsection{Presupuesto Total}
El coste total del proyecto se obtiene sumando los costes de los recursos hardware, software, humanos, licencias y costes indirectos.

\begin{table}[H]
    \centering
    \begin{tabular}{|l|c|c|}
        \hline
        \rowcolor[HTML]{C0C0C0} 
        \textbf{Concepto} & \textbf{Coste Mensual} & \textbf{Coste Total} \\ \hline
        \textbf{Recursos Hardware} & - & 1.000 \euro \\ \hline
        Ordenador portátil & - & 800 \euro \\ \hline
        Teléfono móvil para pruebas & - & 200 \euro \\ \hline
        \textbf{Recursos Humanos} & 1.750\euro & 10.500 \euro \\ \hline
        Programador Junior (6 meses) & 1.750\euro & 10.500 \euro \\ \hline
        \textbf{Licencias} & 52,68\euro & 316,12\euro \\ \hline
        Windows 11 Home License & - & 145 \euro \\ \hline
        Visual Paradigm Standard Edition & 12,23\euro & 73,38\euro \\ \hline
        Figma Professional & 20,45\euro & 122,7\euro \\ \hline
        \textbf{Costes Indirectos} & 262,5\euro & 1.575\euro \\ \hline
        Electricidad & 50\euro & 300\euro \\ \hline
        Internet & 30\euro & 180\euro \\ \hline
        Espacio de trabajo (estimado) & 182,5\euro & 1.095\euro \\ \hline
        \textbf{Coste Total} & 2.065,18\euro & \textbf{13.391,12\euro} \\ \hline
    \end{tabular}
    \caption{Presupuesto detallado del proyecto}
    \label{tab:presupuesto_detallado}
\end{table}




\newpage